\subsection{A Rule for the Season}\label{sec:rule}

The law of the Church asks of every Catholic
in Advent what it asks all year: Mass on Sundays, and in the United States Mass
on 8 December, the solemnity of the Immaculate Conception, a holy day of
obligation. Everything else is offered. But a season kept without any plan
tends to be swallowed by December, and a rule, written down before the First
Sunday and kept modest, is the surest defence.

A short rule might be this. Go to Mass every Sunday and on 8 December, and
listen to the readings as if they were addressed to you. Say the Angelus once a
day. Read the first reading of the day's Mass, at home, every day. Do one work
of mercy each week that costs you something. Go to confession before Christmas
(Section~\ref{sec:confession}).

A fuller rule adds to these. Go to Mass on one or two weekdays as well. Say
Morning or Evening Prayer from the Liturgy of the Hours, or both, and read the
second reading of the Office of Readings, where the Fathers speak. Keep Friday
as a day of penance, as the Church's law asks, and choose for it some fast or
abstinence you can keep without fuss. Set aside money for alms at the start of
the season, before the shopping begins.

If you begin late, begin where you are. The season has a second beginning on 17
December, when the prayers of Mass and Office turn toward Christmas day by
day.

The sick keep Advent from a bed or a chair. Ask the parish to bring Holy
Communion, and ask the priest to come before Christmas if you wish to confess.
Say what you can: the Angelus, a psalm, the day's Gospel, or only the name of
Jesus. Offer the waiting that illness brings, which is a real share in the
waiting of the season.

For those who mourn, the weeks before Christmas can be the hardest of the year.
Advent has something to say to them that no other season says so plainly,
because it is the season of the Lord's return. St Paul wrote his Advent hope to
a grieving church:

\begin{quote}
Then we who are alive, who are left, shall be taken up together with them in
the clouds to meet Christ, into the air: and so shall we be always with the
Lord. Wherefore, comfort ye one another with these words.\par
{\smallskip\noindent\small ---1 Thessalonians 4:16--17 (4:17--18 in the Lectionary's numbering)\par}
\end{quote}

\noindent Pray for your dead by name at the Masses of Advent. The coming the
Church prays for is the one at which they will rise.

Children keep Advent with their hands and eyes. The \emph{Directory on Popular
Piety and the Liturgy} describes the Advent wreath, four candles set on green
branches and lit one more each Sunday until Christmas, and gives it its meaning:
the ages of salvation before Christ came, and the prophets' light growing, one
flame after another, through the long night until the Sun of justice rose
(n.~98). Light it at the family meal and read the Sunday Gospel by it. The same
Directory speaks of the crib set up in the home (n.~104); build it slowly through
the weeks and leave the manger empty until Christmas night. In Hispanic homes and
parishes \emph{Las Posadas} re-enacts Joseph and Mary's search for lodging in
Bethlehem (n.~99), and children can walk in it.

Two older devotions suit a household. The St Andrew prayer is a short greeting of
the hour of the Lord's birth; A. A. Lings, in \emph{Our Favorite Novenas},
says that it is piously believed that whoever says it fifteen times a day, from the feast of St
Andrew on 30 November until Christmas, will obtain what is asked. Its words are
in Appendix~\ref{sec:prayers}.
And St Alphonsus Liguori, in \emph{The Incarnation, Birth and Infancy of Jesus
Christ}, set a Christmas novena of meditations beginning on 16 December. A family can read one each evening of the last nine days. From 17
December the great O antiphons, one for each evening, can be sung or said with
the children before supper, and the oldest child can be given the task of
lighting the candle.
